\documentclass[11pt]{article}
\usepackage{varnernotes}

\newcommand{\E}{\mathbb{E}}

\begin{document}

\lecturenotes{12a}{Stochastic Multi-Armed Bandits}{Fall 2026}{November 10, 2026}{Jeffrey D. Varner}

\learningobjectives
  {Formulate a stochastic bandit}
  {Define arms, reward distributions, policies, action counts, optimal mean reward, and cumulative pseudo-regret.}
  {Balance learning and earning}
  {Compare $\varepsilon$-greedy, upper-confidence-bound, and Thompson-sampling decisions through their uncertainty mechanisms.}
  {Audit financial applications}
  {Choose observable rewards and identify nonstationarity, dependence, costs, and off-policy evaluation as threats to a bandit backtest.}

\section{Learning While Acting}

A sequential decision maker often faces an awkward feedback loop: choosing an action produces both a reward and information about that action. Exploiting the option that currently looks best earns near-term reward; exploring a less-certain option can improve later decisions. The stochastic multi-armed bandit is the smallest mathematical model that isolates this tension.

The name comes from a row of slot machines---historically called one-armed bandits---whose payoff rates are unknown. At round $t=1,2,\ldots,T$, an agent chooses one of $K$ arms,
\begin{equation}
  \label{eq:action}
  A_t\in\mathcal A:=\{1,\ldots,K\},
\end{equation}
and observes only the reward $R_t=R_t(A_t)$ generated by the chosen arm. In the stationary stochastic model, arm $a$ has an unknown, time-invariant reward distribution $\nu_a$ with finite mean
\begin{equation}
  \label{eq:means}
  \mu_a:=\E_{R\sim\nu_a}[R],
  \qquad
  \mu^*:=\max_{a\in\mathcal A}\mu_a.
\end{equation}
Conditional on the selected arms, rewards are usually assumed independent through time. These assumptions are a model, not consequences of calling a problem a bandit.

\begin{figure}[ht]
  \centering
  \begin{tikzpicture}[>=Latex,font=\sffamily\small,node distance=2.1cm]
    \node[draw=vnink,very thick,rounded corners,fill=vnsoft,minimum width=2.8cm,minimum height=1.05cm,align=center] (policy) {policy $\pi$\\history $H_{t-1}$};
    \node[draw=vnink,very thick,rounded corners,fill=white,minimum width=4.2cm,minimum height=1.05cm,align=center,right=of policy] (arms) {unknown arms\\$\nu_1\quad\nu_2\quad\cdots\quad\nu_K$};
    \draw[vnlink,very thick,->] (policy) -- node[above] {choose $A_t$} (arms);
    \draw[vncarnelian,very thick,->,bend left=24] (arms) to node[below,align=center] {observe only\\$R_t(A_t)$} (policy);
    \node[below=1.15cm of policy,align=center,text=vnlink] (explore) {explore\\reduce uncertainty};
    \node[below=1.15cm of arms,align=center,text=vncarnelian] (exploit) {exploit\\use current evidence};
    \draw[thick,vnmuted,<->] (explore) -- node[above,text=vnmuted] {tradeoff} (exploit);
  \end{tikzpicture}
  \caption{Bandit feedback is partial: a round reveals the outcome of the selected arm, not the rewards that the unselected arms would have produced.}
  \label{fig:bandit-loop}
\end{figure}

This is not yet a full reinforcement-learning problem. A basic bandit has no evolving state and no action-dependent state transition. A contextual bandit observes covariates before acting but still treats each reward as immediate. Full reinforcement learning is needed when today's action changes tomorrow's state and opportunity set.

\section{Policies, Evidence, and Regret}

The history available before round $t$ is
\begin{equation}
  \label{eq:history}
  H_{t-1}:=(A_1,R_1,\ldots,A_{t-1},R_{t-1}).
\end{equation}
A policy $\pi$ maps this history, and possibly internal randomization, to a distribution over $A_t$. Define the number of observations and empirical mean for arm $a$ before round $t$ by
\begin{equation}
  \label{eq:estimates}
  N_a(t):=\sum_{s=1}^{t-1}\mathbf 1\{A_s=a\},
  \qquad
  \widehat\mu_a(t):=\frac{1}{N_a(t)}
  \sum_{s=1}^{t-1}R_s\mathbf 1\{A_s=a\},
\end{equation}
when $N_a(t)>0$. A warm start that selects every arm once makes these estimates defined.

Expected cumulative pseudo-regret compares the policy with an oracle that always selects an arm with mean $\mu^*$:
\begin{equation}
  \label{eq:regret}
  \operatorname{Reg}_T(\pi)
  :=T\mu^*-\E_\pi\!\left[\sum_{t=1}^{T}R_t\right].
\end{equation}
It measures the expected opportunity cost of learning under the assumed stationary model. It is not the realized difference between two noisy sample paths, and it does not include trading costs unless those costs are built into $R_t$.

\begin{theorem}{Regret decomposition}{regret-decomposition}
Let $\Delta_a:=\mu^*-\mu_a\geq 0$ and let $N_a(T+1)$ count selections of arm $a$ through round $T$. Under the stationary stochastic model,
\begin{equation}
  \label{eq:regret-decomposition}
  \operatorname{Reg}_T(\pi)
  =\sum_{a=1}^{K}\Delta_a\,\E_\pi[N_a(T+1)].
\end{equation}
Thus regret is controlled by how often a policy selects suboptimal arms and by the mean-reward gaps attached to those selections.
\end{theorem}

To see the result, group expected rewards by arm:
$\E[\sum_t R_t]=\sum_a\mu_a\E[N_a(T+1)]$, then use
$\sum_a N_a(T+1)=T$. The identity also explains why learning is difficult when gaps are small: many observations may be required to distinguish nearly equal means.

\section{Three Ways to Explore}

An $\varepsilon$-greedy policy chooses an empirical-best arm with probability $1-\varepsilon_t$ and otherwise samples an arm at random:
\begin{equation}
  \label{eq:epsilon-greedy}
  A_t=
  \begin{cases}
    \text{a random arm in $\mathcal A$}, & \text{with probability }\varepsilon_t,\\
    \arg\max_a\widehat\mu_a(t), & \text{with probability }1-\varepsilon_t.
  \end{cases}
\end{equation}
It is easy to implement, but uniform exploration ignores which arms are uncertain. A constant $\varepsilon_t=\varepsilon>0$ spends a constant fraction of rounds exploring and therefore generally incurs regret proportional to $T$. A decreasing schedule can reduce this cost, although a poorly chosen schedule may stop exploring before the rankings are reliable.

For rewards bounded in $[0,1]$, the UCB1 rule selects
\begin{equation}
  \label{eq:ucb}
  A_t\in\arg\max_a
  \left\{
    \widehat\mu_a(t)+\sqrt{\frac{2\log t}{N_a(t)}}
  \right\}.
\end{equation}
The first term is estimated value; the second is an uncertainty bonus. An arm becomes attractive because it has performed well or because it has not been sampled enough. Under the bounded, independent, stationary assumptions, this optimism principle gives logarithmic gap-dependent regret rather than the linear regret of constant-$\varepsilon$ exploration.

Thompson sampling maintains a probability model for each arm's unknown parameter. At each round it draws one parameter value from each posterior distribution and chooses the arm with the largest sampled mean. Posterior uncertainty automatically produces exploration; posterior concentration increasingly favors exploitation. In the Bernoulli--Beta case developed in L11b, the update is conjugate and requires only success and failure counts.

\begin{figure}[ht]
  \centering
  \begin{tikzpicture}[x=1.2cm,y=0.92cm,>=Latex,font=\sffamily\small]
    \draw[->,vnmuted] (0,0)--(7.7,0) node[right] {arm};
    \draw[->,vnmuted] (0,0)--(0,4.0) node[above] {mean reward};
    \foreach \x/\m/\u/\lab in {1.2/2.3/0.35/1,3.3/2.65/1.05/2,5.4/3.0/0.45/3}{
      \draw[vnlink,very thick] (\x,\m-\u)--(\x,\m+\u);
      \draw[vnlink,thick] (\x-0.18,\m-\u)--(\x+0.18,\m-\u);
      \draw[vnlink,thick] (\x-0.18,\m+\u)--(\x+0.18,\m+\u);
      \fill[vncarnelian] (\x,\m) circle (2.4pt);
      \node[below] at (\x,0) {$\lab$};
    }
    \draw[vnmuted,dashed] (0,3.7)--(6.7,3.7);
    \node[right,align=left,text=vnmuted] at (5.65,3.66) {largest upper\\confidence bound};
  \end{tikzpicture}
  \caption{Schematic UCB decision. Arm 2 is not the empirical leader, but its wider uncertainty interval can justify another observation. Intervals here are conceptual; \eqnref{ucb} defines the actual index.}
  \label{fig:ucb}
\end{figure}

\begin{remark}{The objective determines the algorithm}{objective}
Regret minimization is a cumulative-reward objective. Best-arm identification instead spends a fixed sampling budget to recommend one arm accurately at the end. Pure exploration, risk constraints, switching costs, delayed outcomes, and nonstationarity produce different problems and different guarantees.
\end{remark}

\section{From Arms to Financial Decisions}

An arm could represent an asset, signal, execution venue, or trading rule. The reward must be observable after the action and aligned with the decision goal. Examples include net excess return over a benchmark, a binary indicator that excess return is positive, or a utility-adjusted outcome. Binary rewards simplify inference but discard the magnitude and tail shape of gains and losses. Gross return is misleading when turnover, spread, fees, financing, and market impact differ across arms.

The textbook model is especially fragile in markets:
\begin{itemize}[leftmargin=1.3em,itemsep=2pt,topsep=2pt]
  \item reward distributions and the identity of the best arm change through regimes;
  \item asset rewards share market and factor shocks, so arms are not independent experiments;
  \item outcomes may be delayed or overlap, and prices can move while an action is executed;
  \item the available asset universe and data history can contain survivorship and look-ahead bias; and
  \item risk limits, inventory, and transaction costs couple decisions through time.
\end{itemize}
Discounted estimates, rolling windows, change-point methods, contextual information, or a stateful control model may be more appropriate, but each changes the question being answered.

Historical data also reveal outcomes for assets that a live bandit would not have selected. Treating this full-information panel as if it were bandit feedback changes the experiment. Conversely, replaying only logged actions can be selection-biased. Off-policy evaluation generally requires known logging propensities and support for every target-policy action; a common inverse-propensity estimate weights an observed reward by the ratio of target to logging action probability. Without those conditions, a backtest cannot recover the counterfactual rewards merely by resampling history.

\notebooklink
  {L12b Preview: Binary Bandit Ticker Picker}
  {https://github.com/varnerlab/CHEME-5660-CourseRepository-Fall-2026/blob/main/lectures/week-12/L12b/CHEME-5660-L12b-Example-BBBP-Ticker-Picker-Fall-2026.ipynb}
  {Preview the Bernoulli--Beta implementation used next lecture, and audit how ticker outcomes are converted into binary rewards.}

\keytakeaways
  {In this lecture, we turned the exploration--exploitation tension into a stochastic bandit model and connected policy choices to expected cumulative regret.}
  {Feedback is partial}
  {A bandit observes the chosen arm's reward, so every action simultaneously earns a payoff and changes the evidence available for later decisions.}
  {Exploration needs structure}
  {$\varepsilon$-greedy randomizes directly, UCB adds an optimism bonus, and Thompson sampling randomizes according to posterior uncertainty.}
  {Financial rewards require an audit}
  {Nonstationarity, common shocks, costs, risk, and counterfactual evaluation can invalidate textbook assumptions even when the algorithm is coded correctly.}
  {Next, specialize the framework to binary ticker outcomes with Bernoulli likelihoods and Beta posteriors.}

\begin{readinglist}
  \item Herbert Robbins, \href{https://doi.org/10.1007/978-1-4612-5110-1_26}{``Some Aspects of the Sequential Design of Experiments,''} \emph{Bulletin of the American Mathematical Society} 58(5), 527--535 (1952).
  \item Tze Leung Lai and Herbert Robbins, \href{https://doi.org/10.1016/0196-8858(85)90002-8}{``Asymptotically Efficient Adaptive Allocation Rules,''} \emph{Advances in Applied Mathematics} 6(1), 4--22 (1985).
  \item Peter Auer, Nicol\`o Cesa-Bianchi, and Paul Fischer, \href{https://doi.org/10.1023/A:1013689704352}{``Finite-Time Analysis of the Multiarmed Bandit Problem,''} \emph{Machine Learning} 47, 235--256 (2002).
  \item Daniel J. Russo et al., \href{https://doi.org/10.1561/2200000070}{``A Tutorial on Thompson Sampling,''} \emph{Foundations and Trends in Machine Learning} 11(1), 1--96 (2018).
\end{readinglist}

\vfill
{\sffamily\footnotesize\color{vnmuted}\textbf{Educational use and risk notice.} These notes are for instruction only. A bandit policy does not remove estimation, market, liquidity, or implementation risk.}

\end{document}
